%----------------------------------------------------------------------------
\chapter{Megvalósítás}
%----------------------------------------------------------------------------

\section{Az alkalmazás architektúrája}
Az alkalmazás fejlesztése során fontos szempont volt egy jól struktúrált, könnyen bővíthető és karbantartható architektúra létrehozása. A rendszert három fő komponensből építettem fel: frontend réteg, backend réteg és AI service réteg. Ezek a komponensek mind egymástól elkülönítve működnek, kommunikációjuk HTTP alapú API hívásokon keresztül történik.



\section{Frontend megvalósítása}

Az alkalmazás kliensoldali része Flutter keretrendszer használatával készült. \cite{flutter} Ez lehetőséget biztosít modern, platformfüggetlen mobil vagy akár web és desktop alkalmazások fejlesztésére. A Flutter előnye nem más, mint hogy egyetlen kódbázisból több platformra is készíthető alkalmazás. Az Echo alkalmazásban a frontend réteg biztosítja a kapcsolatot a felhasználó és a rendszer többi komponense között, ezen keresztül történik minden fontos funkció elérése, mint a regisztráció, beszélgetés, gyakorlás, hibák és szókincs megtekintése.

A projekt felépítésénél első számú szempont az átláthatóság volt, ezért a fájlokat különböző mappákba csoportosítottuk, funkcionális egységek szerint elrendezve. A \textit{lib} mappában három fontos rész különíthető el: \textit{core}, \textit{features}, \textit{shared} mappa. Ez a szerkezett hozzájárult ahhoz, hogy az alkalmazás részei jól elkülönüljenek, ugyanakkor könnyen bővíthetők legyenek újabb funkcionalitásokkal.

A \textit{core} mappa tartalmazza az egész alakalmazás működéséhez szükséges elemeket. Ide tartoznak a konstans értékek, globális providerek, téma beállítások. A hálózati kommunikáció megvalósításához Dio csomagot  \cite{dio} használtam, amellyel HTTP kéréseket tudtunk küldeni a backend felé. Az API végpontok elérhetősége .env fájlban van eltárolva flutter dotenv csomag segítségével. \cite{flutterdotenv}

A \textit{features} mappa tartalmazza az alkalmazás fő funkcióit. Minden feature külön mappába került. Ez a feature alapú felépítés nagyon hasznos, mivel minden funckionalitás saját egységként kezelhető. A fejlesztés során nagyon megkönnyíti így a dologt, mivel könnyebben átlátható az, hogy melyik működéshez melyik fájl tartozik.

Például, a \textit{conversation} modul külön \textit{data} és \textit{presentation} rétegre oszlik. A \textit{data} réteg az API kommunikációhoz és adatmodellekhez kapcsolódó fájlokat tartalmazza. Ezek felelősek a backenddel való kommunikációért, mint például az üzenetek kezeléséért, valamint hangalapú beszélgetések válaszainak feldolgozásáért. A \textit{presentation} réteg viszont a felhasználói felülethez kapcsolódó komponenseket tartalmazza.  

Állapotkezelés szempontjából a \textit{flutter riverpod} \cite{riverpod} csomagot vettem használatba. Ez a csomag hozzájárul ahhoz, hogy az alkalmazás különböző részei könnyen elérjék és frissítsék a szükséges adatokat, például a bejelentkezett felhasználó adatait, a beszélgetések állapotát vagy a betöltési folyamatokat. Ennek az előnye az, hogy jól illeszkedik a moduláris felépítéshez, és hozzájárul az állapotkezeléshez.

Az alkalmazásban az egyik fontos funkcionalitás a hangalapú kommunikáció. Erre három különböző csomagot használtam fel. A \textit{record} segítségével a felhasználó hangot tud rögzíteni, amelyet később továbbíthatják a backend felé, feldolgozásra. Az \textit{audioplayers} csomag a hangválaszok lejátszására szolgálnak, míg a \textit{path provider} a fájlok ideiglenes vagy helyi tárolásában nyújt segítséget. \cite{record,audioplayers,pathprovider}

Mégegy fontos mappa a \textit{shared} mappa. Itt olyan újrafelhasználható elemeket tárolunk, amelyeket több képernyőn is megjeleníthetünk. Ide tartoznak a közös widgetek, amelyekkel egységessé varázsoljuk a felületet. Ennek a mappának segítségével csökkentettem a kódismétlést és megkönnyítette a későbbi módosításokat.

Összességében a frontend megvalósítása során egy jól struktúrált, funkcionalitás alapú Flutter alkalmazás jött létre, amely jól kezelhető, képes kommunikálni megfelelően az APIval, valamint támogatni a szöveges és hangalapú nyelvtanulási folyamatokat.

\newpage

\section{Backend megvalósítása}

Az alkalmazás szerveroldali komponensét ASP.NET Core Web API hasnzálatával készítettem. \cite{aspnetcore} Itt valósul meg a felhasználói hitelesítés, üzleti logika, adatbázissal való kommunikáció, illetve a mesterséges intelligenciát biztosító szolgáltatásokkal való kapcsolat fenntartása. A fejlesztés során fontos szempont volt a kód átláthatósága, karbantarthatósága, így esett a választás a Clean architektúrára. \cite{martin2017cleanarchitecture}

A projekt négy fő rétegre bontható: API, Application, Core és Infrastructure. Ezek a rétegek mind elkülönült felelősséggel rendelkeznek, kevés függőséggel, így egyszerű, hosszútávú fejlesztéséhez járul hozzá.

\begin{itemize}
    \item \textbf{API réteg:}
    
    Az API réteg felelős a kliens alkalmazás és a backend közötti kapcsolat létesítésében. Itt HTTP kérések segítségével kommunikálhatnak egymással.

    Különböző mappák találhatóak ebben a rétegben, mindegyiknek különböző, de nagyon fontos szerepe van a projekben.

    A Controller mappában különböző végpontok tárolódnak el. Ez a mappa ezeknek a végpontok kezeléséért, tárolásáért, rendszerezéséért felelős. Többek között ezek a végpontok biztosítják a felhasználói hitelesítést, beszélgetés kezelést, beállításokat, és még sok egyebet.

    A DTO mappában szerepelnek azok az osztályok, amelyek az adatátvitelt szolgálják a kliens és szerver között. Ezek az osztályok nagyon fontos szerepet játszanak a projektben, mivel hozzájárulnak ahhoz, hogy csak a szükséges adatokat továbbítsuk a kliens számára. Ezáltal növeljük a rendszer biztonságát és rugalmasságát.

    A Mappings mappa tartalmazza az AutoMapper  \cite{automapper} konfigurációit. Itt automatikusan átalakítjuk az entitásokat DTO objektumokká és ugyanúgy fordítva. Ez olyan szempontból hasznos, hogy leredukálja a manuális konvertálások számát, egyszerűsíti a fejlesztést, így időt spórol ugyanakkor van egy központi helye ezeknek az átalakításoknak.

    Van egy Middleware \cite{aspnetcoreMiddleware} is a rendszerben, amelyben a globális hibakezelést valósítottam meg. Ennek feladata, hogy központilag kezelje az előforduló hibákat és egységes formátumú válaszokat küldjön a kliens számára.

    \item \textbf{Application réteg}

    Ez a réteg tartalmazza az alkalmazás üzleti logikáját. Itt olyan szolgáltatásokkal találkozhatunk, amelyek a különböző folyamatok működését valósítják meg.

    Itt történik például a felhasználók kezelése, hibák eltárolása, AI szolgáltatásokkal kapcsolatos implementációk. Itt nem szabad adatbázis specifikus megvalósításokat tárolni, kizárólag üzleti logikát, és interfészek segítségével valósul meg a kommunikáció a külső komponensekkel.

    \item \textbf{Core réteg}

    Ez az a réteg, amely a rendszer legbelső része. Itt vannak az alkalmazás alapvető elemei, domain logikái.

    Az Entities mappa az adatbázis entitásait tárolja, ezek az üzleti objektumok. Ezeknek segítségével és az Entity Framework segítségével jöttek létre code first módon az adatbázis táblái.

    Az Enums mappában tároljuk a felsorolásokat, ennek segítségével egységes kezelési módot alakítunk ki az különböző állapotok és kategóriák között.

    Az Interfaces mappa egy nagyon fontos hely, mivel ezek határozzák meg a szolgáltatások működését. A dependency injection ennek segítségével valósul meg, amely lehetővé teszi azt implementációk egyszerű cseréjét és a rendszer jobb tesztelhetőségét. \cite{dotnetDI}

    \item \textbf{Infrastructure réteg}

    Itt a külső szolgáltatások és adatbázis eléréseket valósítom meg.

    A Data mappa szolgál az adatbázis kapcsolat konfigurációjában, a DbContext osztály pedig a táblák beállításait és a kapcsolatok kezeléseit biztosítja.

    A Repositories mappa olyan osztályokat tartalmaz, amelyek az adatbázisműveletek implementációit tartalmazza. Ezek az osztályok lekérdezéséért, módosításáért, létrehozásáért és törléséért felelősek. A repository minta lett itt alkalmazva, amely elkülöníti az adatkezelést az üzleti logikától.

    A Services mappában olyan szolgáltatások vannak, amelyek külső rendszerekkel kommunikálnak. Itt leginkább a repository osztályokat hívjuk, illetve itt valósul meg az AI szolgáltatások hívása is.

    Nem utolsó sorban a Utils mappát is megtalálhatjuk itt, amely különböző segédosztályokat tartalmaz, amelyeket bárhol a programban felhasználhatunk.

    \item \textbf{Program konfiguráció}

    A rendszer indulásakor a Program.cs fájl központi szerepet játszik. Ez végzi el az alkalmazás konfigurációját. Itt állítódik be az adatbázis kapcsolat, dependency injection konfigurálása, jwt hitelesítés beállítása, middlewarek regisztrációja. \cite{jwt}
\end{itemize}

% \newpage

\section{AI service megvalósítása}

A mesterséges intelligenciához kapcsolodó funkciók külön Python alapú mikroszolgáltatásokként lettek megvalósítva. \cite{fowlerMicroservices} Ez a felépítés nagyon előnyös, mivel így elkülönül a .NET backend az AI funkcióktól. Ezzel külön fejleszthetővé válik minden rész, szükség esetén könnyen módosítani lehet minden részt, vagy akár cserélni is. 

Az AI réteg három fő szolgáltatásból áll össze: beszédfelismeréshez a Whisper service járul hozzá, a nyelni modell a Phi3 service , a szövegfelolvasást pedig a Piper service végzi. 

Mindhárom szerviz hasonló projektstruktúrával rendelkezik. Mindhárom mappában külön mappa van az API útvonalaknak, konfigurációs fájloknak, adatmodelleknek és a logikát tartalmazó osztályoknak. Az api mappa endpointokat tartalmaz, a core konfigurációkat, models a kérés és válaszmodelleket, és a service az adott AI funkció tényleges megvalósítását.

Elsőszámú AI service a Whisper, amely azért felelős, hogy a beszédet szöveggé alakítsa. Itt faster whispert használtam, amely egy gyors változata a modellnek. \cite{whisper,fasterwhisper} Az alkalmazásban ez lesz az a komponens, ami elsőként hívódik meg, amikor a felhasználó a mesterséges intelligenciához beszél, ez pedig szöveges formátummá fogja átalakítani a következő lépéshez.

A Phi3 service a nyelvi intelligenciát biztosítja. Ez beszélgetésre és utasításkövetésre optimalizált nyelvi modell. Feladata az, hogy a megadott szövegre utasítás szerinti válaszokat generáljon, felismerje és kijavítsa a nyelvi hibákat a visszaadott szövegben. \cite{phi3}

A Piper service a harmadik fázishoz járul hozzá a válaszadásban, amikor a szöveges választ, amelyet a Phi3 generált, beszéddé alakítja. Ez azért fontos, mert egy valós beszélgetést szerettem volna leszimulálni, hogy a felhasználó könnyebben memorizálja és hozzászokjon, milyen a valódi kommunikáció. \cite{piper}

A backend API HTTP kéréseken keresztül kommunikál ezekkel a mikroszolgáltatásokkal.